\chapter{Introducción}

\lettrine{E}{l} desarrollo de entornos virtuales contemporáneos, como videojuegos y simulaciones interactivas, exige dotar a sus agentes de la inteligencia suficiente para interactuar coherentemente con su entorno.
El comportamiento de estas entidades es clave para garantizar la credibilidad, funcionalidad y realismo del entorno virtual, especialmente en escenarios complejos y dinámicos.
A pesar de que se trate de un problema que lleva décadas siendo explorado en el campo de la inteligencia artificial~\cite{russell1995modern}, a día de hoy, dista de ser trivial cuando se persigue un comportamiento robusto y eficiente.

Dentro de este ámbito, un requisito fundamental es la \textbf{navegación autónoma}.
Este proceso permite a un agente desplazarse desde un punto de origen hasta un destino específico calculando una ruta viable y esquivando, de manera eficiente, la geometría y los obstáculos presentes en la escena.
En el contexto de videojuegos, este comportamiento recae habitualmente en los personajes no jugables, conocidos en inglés como \acrfull{npc}.
Su capacidad para desenvolverse de forma autónoma condiciona en gran medida la percepción de inteligencia del sistema. 
Abordar este reto no solo implica determinar el camino más corto, sino también garantizar que el movimiento resultante sea compatible con las restricciones físicas del entorno y con el comportamiento esperado del agente~\cite{reynolds1999steering}.

El problema de la búsqueda de rutas, conocido en la literatura como \textit{pathfinding}, constituye uno de los casos de estudio más explorados en el campo de la inteligencia artificial aplicada a entornos virtuales~\cite{Millington}.
En la actualidad, se considera un problema ampliamente resuelto para agentes terrestres en escenarios estáticos~\cite{abd2015comprehensive}, hasta el punto de que los principales motores de desarrollo, tanto comerciales como de código abierto, integran soluciones nativas y robustas para este fin.
Ejemplos representativos de ello son el Navigation System en Unreal Engine~\cite{UnrealNavSystem}, el sistema nativo de NavMesh en Unity~\cite{UnityNavMesh} o el componente NavigationServer3D en Godot~\cite{GodotNavServer}, los cuales permiten implementar comportamientos de navegación complejos con un coste computacional asumible.

La gran mayoría de estas herramientas se fundamentan en la biblioteca de código abierto Recast Navigation~\cite{RecastNavigation}.
Este sistema permite generar una \textbf{malla de navegación}, que recibe el nombre de NavMesh, y que representa una \textbf{aproximación 2.5D} del entorno transitable.
Aunque el escenario se modela en tres dimensiones, el cálculo de rutas se proyecta sobre superficies bidimensionales que definen las zonas donde los agentes pueden apoyarse o caminar, tal y como puede apreciarse en la Figura~\ref{fig:navmesh25d}.
A partir de esta geometría, se construye un grafo sobre el cual se aplican algoritmos de búsqueda heurística, siendo el algoritmo \gls{astar} el caso estándar de aplicación~\cite{astar}.

\begin{figure}[htp]
  \centering
  \includegraphics[width=0.75\textwidth]{imaxes/navmesh25d.png}
  \caption[Ejemplo de malla 2.5D generada por NavMesh sobre un entorno tridimensional.]
  {Ejemplo de malla 2.5D generada por NavMesh sobre un entorno tridimensional.
  Imagen extraída de~\cite{UnityNavMeshAgentManual}.}
  \label{fig:navmesh25d}
\end{figure}

\section{Motivación}

La necesidad de crear una herramienta que dote a agentes voladores de la capacidad de moverse a través de un espacio tridimensional surge del limitado soporte que ofrecen las tecnologías dominantes en el desarrollo de videojuegos actual.
Las soluciones tradicionales basadas en mallas 2.5D resultan ideales para personajes terrestres, ya que presuponen que el agente está anclado al suelo y que su movimiento puede reducirse a un plano horizontal cuya altura es una propiedad derivada del terreno.
Sin embargo, este supuesto no se cumple para entidades con capacidad de vuelo, tales como drones, naves espaciales, enjambres de insectos o criaturas submarinas como peces, las cuales requieren operar con \textbf{seis grados de libertad}: tres traslaciones a lo largo de los ejes $X$, $Y$ y $Z$ y tres rotaciones independientes alrededor de esos mismos ejes, conocidas como \textit{pitch}, \textit{yaw} y \textit{roll}, tal y como se ilustra en la Figura~\ref{fig:grados_libertad}.
Ante la posibilidad de que los agentes puedan modificar su trayectoria en cualquier dirección y altura, un plano transitable resulta insuficiente y se hace imprescindible un \textbf{volumen de navegación}.
Este debe representar la \textbf{totalidad del espacio libre} del entorno y no solo las superficies, logrando garantizar que ninguna ruta viable o alternativa óptima quede excluida durante el cálculo de trayectorias.

\begin{figure}[htp]
  \centering
  \includegraphics[width=0.75\textwidth]{imaxes/6dof.png}
  \caption[Visualización de los 6 grados de libertad.]
  {Visualización de los 6 grados de libertad.
  Imagen extraída de~\cite{6DoF_IIA}.}
  \label{fig:grados_libertad}
\end{figure}

Por consiguiente, el salto hacia una navegación puramente volumétrica plantea un desafío técnico complejo.
Al añadir la altura como una dimensión libre, la gran cantidad de espacio a analizar dispara el coste computacional tanto en tiempo como en memoria~\cite{lavalleplanning}.
Si esta complejidad no se gestiona mediante estructuras de datos espaciales altamente especializadas como \glspl{octree}~\cite{octree} o representaciones jerárquicas equivalentes, el cálculo de trayectorias provocará un consumo crítico de memoria y una degradación severa en la tasa de \acrfull{fps}, comprometiendo la viabilidad comercial del producto.

Ante la falta de un soporte oficial eficiente en los motores dominantes, los estudios de desarrollo se ven obligados a implementar soluciones \gls{adhoc} para cada proyecto, o bien a intentar adaptar de forma artificial las tecnologías 2.5D existentes~\cite{Brewer2015NavMesh}.
Ambas alternativas suelen derivar en un incremento significativo de los costes de producción, así como en una pérdida de tiempo que podría dedicarse al diseño del propio videojuego.
Surge, por lo tanto, la necesidad evidente de diseñar un sistema \textbf{genérico, optimizado y reutilizable} que resuelva el problema de la navegación tridimensional de manera estructural.

\section{Objetivos}
\label{sec:objetivos}

El objetivo general de este Trabajo de Fin de Grado consiste en el \textbf{diseño, implementación y validación en el motor Unity de un sistema de navegación volumétrica eficiente y optimizado}.
El resultado final tendrá un enfoque de producto: un paquete de herramientas modular, de código abierto, listo para su integración inmediata en proyectos y entornos de producción reales.

A partir de este objetivo general, se establecen los siguientes objetivos específicos:

\begin{itemize}
    \item \textbf{Diseño arquitectónico:} Diseñar una arquitectura modular y escalable que garantice la eficiencia computacional y facilite la futura integración de nuevas funcionalidades.
    \item \textbf{Representación eficiente del espacio:} Implementar una estructura de datos espacial que represente el volumen navegable de forma compacta y optimizada para consultas en tiempo real.
    \item \textbf{Sistema de serialización:} Desarrollar un módulo de almacenamiento en disco para guardar la representación generada, evitando el coste de su reconstrucción en tiempo de ejecución.
    \item \textbf{Búsqueda de rutas:} Adaptar un algoritmo de búsqueda heurística tridimensional para hallar trayectorias viables dentro del espacio libre del volumen.
    \item \textbf{Postprocesado de trayectorias:} Desarrollar un algoritmo de suavizado de caminos que optimice la trayectoria resultante mediante la eliminación de giros bruscos y nodos redundantes.
    \item \textbf{Integración en Unity:} Desarrollar una herramienta nativa para el editor del motor que simplifique la configuración del sistema y su visualización mediante Gizmos y paneles personalizados.
    \item \textbf{Accesibilidad y usabilidad:} Proveer una interfaz intuitiva que permita la configuración del sistema tanto a perfiles técnicos como a diseñadores de niveles y artistas.
    \item \textbf{Análisis de rendimiento:} Evaluar cuantitativamente la eficiencia global del sistema y de sus componentes críticos bajo diferentes escenarios y cargas de prueba.
    \item \textbf{Distribución modular:} Empaquetar el sistema final bajo los estándares oficiales de Unity para facilitar su integración en otros proyectos y su potencial publicación en la Unity Asset Store.
\end{itemize}

\section{Estructura de la memoria}

La presente memoria se organiza en los siguientes capítulos:

\begin{itemize}
    \item \textbf{Capítulo 1: Introducción.} El presente capítulo.
    \item \textbf{Capítulo 2: Fundamentos.} Explica los distintos términos y técnicas requeridas para comprender el proyecto.
    \item \textbf{Capítulo 3: Estado de la cuestión.} Examina las soluciones existentes en la literatura y la industria, incluyendo un análisis de \acrfull{dafo}.
    \item \textbf{Capítulo 4: Metodología y planificación.} Detalla el modelo de ciclo de vida de software adoptado, así como la planificación temporal y los recursos empleados.
    \item \textbf{Capítulo 5: Desarrollo.} Profundiza en las distintas fases involucradas en la implementación del sistema.
    \item \textbf{Capítulo 6: Conclusiones.} Sintetiza los resultados obtenidos y evalúa las aportaciones del sistema al desarrollo de videojuegos.
    \item \textbf{Capítulo 7: Trabajos futuros.} Propone extensiones y nuevas líneas de investigación que pueden desarrollarse a partir del sistema implementado.
\end{itemize}
